\documentclass[12pt,a4paper]{article}
\usepackage[utf8]{inputenc}
\usepackage[T1]{fontenc}
\usepackage[ngerman,latin]{babel}
\usepackage{amsmath,amssymb,amsthm}
\usepackage{geometry}
\usepackage{titlesec}
\usepackage{fancyhdr}
\usepackage{microtype}
\usepackage{longtable}

\geometry{
  a4paper,
  left=2.5cm,
  right=2.5cm,
  top=2.5cm,
  bottom=2.5cm
}

\pagestyle{fancy}
\fancyhf{}
\fancyhead[C]{\small Carl Friedrich Gau\ss{} Werke --- Band 4}
\fancyfoot[C]{\thepage}
\renewcommand{\headrulewidth}{0.4pt}

\title{Carl Friedrich Gau\ss{} Werke --- Band 4, Teil 5}
\author{Carl Friedrich Gau\ss}
\date{}

\begin{document}
\maketitle
\thispagestyle{empty}


\section*{Vorrede}
\addcontentsline{toc}{section}{Vorrede}

Der Verfasser dieser Abhandlung hat die zweimalige Wahl der Aufgabe, die
ihren Gegenstand ausmacht, als einen Beweis von der Wichtigkeit betrachten zu
m"ussen geglaubt, welche die k"onigliche Societ"at derselben beilegt, und ist da-
durch aufgemuntert worden, dieser seine schon vor l"angerer Zeit gefundene Auf-
l"osung vorzulegen, wovon ihn sonst die sp"ate von der Preisfrage erhaltene Kennt-
niss abgehalten haben w"urde. Er bedauert, dass der letztere Umstand ihn ge-
n"othigt hat, sich fast nur auf das Wesentliche und auf die Andeutung einiger n"a-
her liegenden Benutzungen f"ur Kartenprojectionen und f"ur die h"ohere Geod"asie
zu beschr"anken, da er ohne die N"ahe des Schlusstermins gern die Entwicklung
einiger Nebenumst"ande noch weiter verfolgt, und die vielseitigen Anwendungen
in der h"oheren Geod"asie ausf"uhrlich bearbeitet haben w"urde, welches er sich nun
f"ur eine andere Zeit und f"ur einen andern Ort vorbehalten muss.

\begin{flushright}
Im December 1822.
\end{flushright}

\newpage

\section*{Allgemeine Aufl"osung der Aufgabe: die Theile einer gegebenen Fl"ache auf einer andern gegebenen Fl"ache so abzubilden, dass die Abbildung dem Abgebildeten in den kleinsten Theilen "ahnlich wird}
\addcontentsline{toc}{section}{Allgemeine Aufl"osung der Aufgabe}

ALLGEMEINE AUFLÖSUNG DER AUFGABE
DIE T HEILE EINER GEGEBENEN FLÄCHE
AUF EINER ANDERN GEGEBENEN FLÄCHE SO ABZUBILDEN
DASS DIE ABBILDUNG DEM ABGEBILDETEN ■ •
IN DEN KLEINSTEN THEILEN ÄHNLICH WIRD.
1.
Die Natur einer krummen Fläche wird durch eine Gleichung zwischen den
sich auf jeden Punkt derselben beziehenden Coordinaten x,y, z bestimmt. Ver-
möge dieser Gleichung kann jede dieser drei veränderlichen Grössen wie eine
Function der beiden andern betrachtet werden. Noch allgemeiner ist es , noch
zwei neue veränderliche Grössen t, u einzuführen, und jede der x, y, z als eine
Function von t und u darzustellen, wodurch, wenigstens allgemein zu reden,
bestimmte Werthe von t und u allemal einem bestimmten Punkte der Oberfläche
angehören , und umgekehrt.
2.
In Beziehung auf eine zweite krumme Fläche sollen X, F, Z, T, TJ ähn-
liche Bedeutungen haben, wie resp. x, y, z, t, u in Beziehung auf die erstere.
3.
Die erste Fläche auf der zweiten abbilden heisst. ein Gesetz festsetzen, nach
welchem einem jeden Punkte der ersten Fläche ein bestimmter Punkt der zwei-
ten entsprechen soll. Dieses wird dadurch geschehen, dass T und U bestimm-
ten Functionen der zwei veränderhchen Grössen t und u gleich 3 g3e setzt werden.

Insofern die Abbildung gewissen Bedingungen Genüge leisten soll, werden diese
Functionen nicht mehr willkürlich sein dürfen. Indem dadurch auch X, Y, Z
zu Functionen von t und u werden, müssen diese Functionen, neben der Bedin-
gung, welche die Natur der zweiten Fläche vorschreibt, auch noch derjenigen
Genüge leisten, welche in der Abbildung erfüllt werden soll.
4.
Die Aufgabe der königlichen Gesellschaft der Wissenschaften schreibt vor,
dass die Abbildung dem Abgebildeten in den kleinsten Theilen ähnlich sein soll.
Es kommt zuvörderst darauf an , diese Bedingung analytisch auszudrücken.
Aus der Differentiation der Functionen von t, u , durch welche <2?,y, z, X , Y, Z
ausgedrückt werden , mögen folgende Gleichungen hervorgehen :
dx = adt -{-ad.u
dy =zhdt -\-h'diu
.' ' d^ = cdt -f-cdw
&X = Adt^Adiu
dF= jBd^+J5'dw
dZ = Cdt-\-C'du
Die vorgeschriebene Bedingung erfordert, erstlich, dass alle von Einem
Punkte der ersten Fläche ausgehende und in ihr liegende unendlich kleine Linien
den ihnen entsprechenden Linien der zweiten Fläche proportional sind, und zwei-
tens, dass jene unter sich dieselben Winkel machen, wie diese.
Ein solches Linear -Element auf der ersten Fläche wird
= \/({aa-i-bb-]-cc)df-}-2{aa-{-bb'-}-cc)dt.du-{-{aa^b'b'-^cc)du^)
und das entsprechende auf der zweiten Fläche
= \/({ÄA-}-BB-{- CC)df-\-2{AÄ-^BB'-\-CC')dt.du-\-(A'Ä'-^B'JB'-^C'C')du^)
Sollen beide, unabhängig von d^ und dw, in einem bestimmten Verhältniss zu
einander stehen , so müssen offenbar die drei Grössen
aa-{-bb-\-cc, aa-{-bb'-\-cc, da-\-b'b'-\-c'c -^
respective den drei folgenden proportional sein :

DIE THEILE EINER GEGEBNEN FLÄCHE AUF EINER ANDERN ABZUBILDEN ETC. 195
ÄÄ-{-BB+CC, ÄÄ'+BB'-j-CC, A'Ä'-^B'B'-\-C'C'
Wenn den Endpunkten eines zweiten Elements auf der ersten Fläche die Werthe
t, u und t-{-^t, u-{-^u
entsprechen , so ist der Cosinus des Winkels , welchen dasselbe mit dem ersten
Elemente macht,
{adi + a'du){ah t -\- a'hu) + (bdt + b'du){bht + b'htt) + {cdt + cdu)(cht + c'hu)
^{{adt + a'dti)*^{bdt + b'du)'' + {cdt-^c'dtiy).{{aot + a'huy+{blt + b'huy + {clt + c'luy)
und für den Cosinus des Winkels zwischen den correspondirenden Elementen
auf der zweiten Fläche ergibt sich ein ganz ähnlicher Ausdruck, wenn nur
a, 6, c, a\ b', c in A, B, C, Ä, B', C verwandelt werden. Offenbar werden beide
Ausdrücke einander gleich, wenn die obige Proportionalität Statt findet, und die
zweite Bedingung wird daher schon mit in der ersten begriffen , welches auch bei
einigem Nachdenken von selbst klar ist.
Der analytische Ausdruck der Bedingung unserer Aufgabe ist demnach, dass
AA-\-BB+CC AA'-\-BB'+CC' A'A'+ B'B'+ C'C
aa-{- bb -\- cc a a' -\- bb' -\- c c' a'a' -{- b' b' -\- c'c'
werden muss, welches eine endliche Function von t und u sein wird, die wir
= mm setzen wollen. Es drückt dann m das Verhältniss aus, in welchem die
Lineargrössen auf der ersten Fläche in ihrer Abbildung auf der zweiten vergrössert
oder verkleinert werden (je nachdem m grösser oder kleiner ist als 1). Dieses
Verhältniss wird, allgemein zu reden, nach den Stellen verschieden sein: in dem
speciellen Falle , wo m constant ist , wird eine vollkommene Aehnlichkeit auch
in den endlichen Theilen , und wenn überdiess m = 1 ist , wird eine vollkom-
mene Gleichheit Statt finden , und die eine Fläche sich auf die andere abwickeln
lassen.
Indem wir Kürze halber
{aa-\-bb-{-cc)df-{-2[aa-{-hb'-{-cc)dt.du-\-{aa-{-b'b'-^cc)du~ = w
setzen , b emerken wir , dass die Differentialgleichung w = 0 zwei Integrationen
zulassen wird. Indem man nemlich das Trinomium w in zwei, in Beziehung auf

dt und du lineare, Factoren zerlegt, miiss entweder der eine oder der andere
Factor = 0 werden, welches zwei verschiedene Integrationen geben wird. Die
eine Integration wird der Gleichung
0 = (^aa-\-bb-\-cc)dt
-\- {a a-\- h h'-{- c c'-\- i \l([a a-{-bh-{-cc) [dd-\- yh'-\- c'c ) — [a d-{- b b'-}- ccf)\du
entsprechen (wo i Kürze halber für \J— l geschrieben ist, indem man sich leicht
überzeugt , dass der irrationale Theil des Ausdrucks imaginär werden muss) ; die
andere einer ganz ähnlichen Gleichung, wenn nur i mit — i vertauscht wird.
Ist also das Integral der erstem Gleichung dieses :
p-\-iq = Const.
wo p und q reelle Functionen von t und w bedeuten, so wird das andere Integral
p — iq = Const.
und die Natur der Sache wird es mit sich bringen , dass
{dp-{-idq) .{dp — idq) oder dp^-\-dq^
ein Factor von to , oder
u) = n(dp^-\-dq^)
werden muss, wo n eine endliche Function von t und u sein wird.
Wir wollen nun das Trinomium, in welches
dX^ + dF=^+dZ'-
übergeht, wenn für dX, dF, dZ ihre Werthe durch T, U, dT, dU substituirt
werden, durch Q bezeichnen, und annehmen, dass auf ähnliche Weise, wie vor-
her ,d ie beiden Integrale der Gleichung Q = 0 diese seien : r
P-{-iQ = Const.
P — iQ = Const.
und
Q = iV(dP^+dQ2)
wo P, Q, N reelle Functionen von T und U bedeuten werden.

DIE THEILE EINER GEGEBNEN FLÄCHE AUF EINER ANDERN ABZUBILDEN ETC. 197
Diese Integrationen lassen sich (die allgemeinen Schwierigkeiten des Inte-
grirens bei Seite gesetzt) offenbar vor der Auflösung unserer Hauptaufgabe aus-
führen.
Wenn nun für T, U solche Functionen von t, u substituirt werden, wo-
bei die Bedingung unsrer Hauptaufgabe erfüllt wird, so geht Q in mmia über,
und es wird -
(dP + idQ).(dP-zdQ) __ mmn
{dp-{-idq).{dp — idq) N
Man sieht aber leicht , dass der Zähler im ersten Theile dieser Gleichung durch
den Nenner nur dann theilbar sein kann, wenn
entweder dP-i-idQ durch dp-{-idq, und dP — idQ durch dp — idq,
oder dP-j-^dQ durch dp — idq, und dP — idQ durch dp-{-idq
theilbar ist. Im ersteren Falle wird demnach dP-\-idQ verschwinden, wenn
^p-i^i^q n= 0. oder P-{-iQ wird constant werden, wenn p-[-iq constant an-
genommen wird, d. i. P-\-iQ wird bloss Function von p-\-iq sein, und eben
so P — iQ Function von p — iq. Im andern Falle wird P-{-iQ Function von
p — iq^ und P — iQ Function von p-{-iq sein. Es ist leicht einzusehen, dass
diese Folgerungen auch umgekehrt gelten, nemlich dass, wenn für P-\-iQ, P — iQ
Functionen von p-\-iq, p — iq (entweder respective, oder verkehrt) angenommen
werden , die endliche Theilbarkeit des ß durch w , und sonach die oben erfor-
derlich gefundene Proportionalität Statt haben wird.
Man überzeugt sich übrigens leicht, dass wenn z.B.
P+iQ=f[p-^iq)
P-iQ=f{p-iq) ,
gesetzt werden , die Beschaffenheit der Function /' schon durch die von / be-
dingt wird. Wenn nemlich unter den constanten Grössen, welche letztere etwa
involviren mag, keine andere als reelle befindlich sind, so wird die andere /"
mit der / ganz identisch sein müssen , damit jedesmal reellen Werthen von p, q
reelle AVerthe von P, Q entsprechen; im entgegengesetzten Falle wird sich /'
von f nur dadurch unterscheiden, dass in den imaginären Elementen von f statt
i überall das entgegengesetzte — i gesetzt werden muss.
Man hat hiernächst

198 ^ ALLGEMEINE AUFLÖSUNG DER AUFGABE
P=\f{:p+iq)-\-^f\v-iq)
i Q = iÄP + iq) - ifip - iq)
oder, was dasselbe ist, indem die Function / ganz willkürlich angenommen wird
(nach Gefallen mit Inbegriff constanter imaginärer Elemente) , wird P dem reel-
len und iQ (bei der zweiten Auflösung — iQ) dem imaginären Theile von
f(p-\-iq) gleich gesetzt, und hieraus sodann vermittelst der Elimination T und
U in der Gestalt von Functionen von t und u dargestellt werden. Hiedurch ist
die vorgegebene Aufgabe ganz allgemein und vollständig aufgelöst.
6.
Wenn p'-\-%q eine beliebige bestimmte Function von p-\-iq vorstellt (in-
dem p\ q reelle Functionen von p, q sind), so sieht man leicht, dass auch
p'-}-«5'' = Const. und p — iq = Const. '
die Integrale der Differentialgleichung w = 0 darstellen; in der That werden
jene mit den obigen
p-\-iq = Const. und p — iq = Const.
resp. ganz gleichbedeutend sein. Eben so werden die Integrale der Differential-
gleichung Q= 0
P'-fiQ'= Const. und P'— iQ' = Const
mit den obigen
P-|-iQ r= Const. und P — iQ = Const.
ganz gleichbedeutend sein, wenn P'-f-iQ' eine beliebige bestimmte Function
von P-\-iQ vorstellt (indem P', Q' reelle Functionen von P, Q sind). Es
erhellet hieraus , d ass in der allgemeinen Auflösung unsrer Aufgabe , welche wir
im vorhergehenden Artikel gegeben haben, auch p, q die Stelle von p, q; und
P', Q' die Stelle von P, Q resp. vertreten können. Wenn gleich die Allgemein-
heit der Auflösung durch eine solche Abänderung nichts gewinnt , so kann doch
zuweilen für die Anwendung eine Form zu diesem , die andere zu jenem Zweck
bequemer sein.

DIE THEILE EINER GEGEBNEN FLÄCHE AUF EINER ANDERN ABZUBILDEN ETC. 199
7.
Wenn die Functionen, welche aus der Differentiation der willkürlichen
Functionen /,/' entspringen, durch cp und cp' r esp. bezeichnet werden, so dass
d ./u = cpu . d u , d ./'u = cp'u . d u, so wird in Folge unsrer allgemeinen Auflösung
dP + idQ f ... dP — idQ ,, . ,
also d^T7d7 = ^(i^ + *^)' d^-tdg = '^ ^ v-^i)
Das Vergrösserungsverhältniss bestimmt sich daher durch die Formel
V 4 — „T" * d ~PM^d-Q~^ • ? (p+*?) • " i'ip—iq)
Wir wollen nun noch unsre allgemeine Auflösung mit einigen Beispielen er-
läutern ,w odurch sowohl die Art der Anwendung , als die Beschaffenheit einiger
dabei noch in Betracht kommenden Umstände am besten ins Licht gesetzt wer-
den wird.
Es seien zuvörderst beide Flächen Ebnen , wo wir
X=T, Y= U, Z= 0
werden setzen können. Die Differentialgleichung ' ' :
(o =: df-\-du^ == 0
gibt hier die beiden Integrale
t-]-iu = Const. , t — tu ^:= Const.
und eben so sind die beiden Integrale der Gleichung Q = d T^-{-dU^ = 0, fol-
gende :
T-\-iU^ Const. , T—iU= Const.
Die beiden allgemeinen Auflösungen der Aufgabe sind demnach:
I. T-^iU = f[t + iu), T-iU=^f\t-iu)
II. T-{-iU = f{t — ii(), T—iU = f{t-{-iu)

Dieses Resultat lässt sich auch so ausdrücken: Indem die Charakteristik / eine
beliebige Function bedeutet, hat man den reellen Theil von f{oo-\-iy) für X,
und den imaginären Theil, mit Weglassung des Factors i, entweder für Y oder
für — Y anzunehmen.
Gebraucht man die Charakteristiken cp, 9' in der Bedeutung des Art. 7
und setzt
cp(.2?4-^» = ^ + ^T], ^'[x — iy) = l — i^
wo offenbar ^ und t] reelle Functionen von oc und y sein werden , so hat man,
in der ersten Auflösung , . .
dX + idF= (^ + ^rj)(d^'-f-^dy)
dX— /dF= (e — fT])(d^ — zdy)
und folglich
dX = IdiOß — T]d^
dF = 'r]d<2?+^dj/
Macht man nun
^ = a.cosy, Tj = a.siny
d<27 = d^ . c os_^ , dy = d5 . sin^
dX= d^S.cosG^, dF= diS.siuG
so dass d^ ein Linearelement in der ersten Ebne , g dessen Neigung gegen die
Abscissenlinie , 6.S das correspondirende Linearelement in der zweiten Ebne
und G dessen Neigung gegen die Abscissenlinie bedeutet , so geben die obigen
Gleichungen
d>S. cos (r ^= a . d 5 . c os(^ -f- 7)
dS^.sin G = a.d^.sin (^-|-7)
und folglich, wenn man. was erlaubt ist, a als positiv betrachtet.
d>S'=:a.dÄ, G = ci-\-y
Man sieht also (in Uebereinstimmung mit Art. 7), dass a das Verhältniss der Ver-
grösserung des Elements ds in der Darstellung dS vorstellt, und. wie gehörig,
von y unabhängig ist; und eben so zeigt die Unabhängigkeit des Winkels 7 von
y, dass alle von einem Punkte ausgehende Linearelemente in der ersten Ebne

DIE THEILE EINER GEGEBNEN FLÄCHE AUF EINER ANDERN ABZUBILDEN ETC. 201
durch Elemente in der zweiten Ebne dargestellt werden, die unter sich und, wie
wir hinzufügen können, in demselben Sinn, dieselben Winkel bilden, wie jene.
Wählt man für f eine linearische Function, so dass /u = A-^Bu, wo
die Constanten Coefficienten von der Form sind
Ä = a-\-bi, B ^= c-\-ei
so wird r^M =: B = c-\-ei
also Q = \J [ cc-\-€e), y = Arc.tang-^
Das Vergrösserungsverhältniss ist folglich in allen Punkten constant, und die
Darstellung dem Dargestellten durchaus ähnlich.
Für jede andere Function f wird (wie man leicht beweisen kann)" das Ver-
grösserungsverhältniss nicht c onstant sein , und die Aehnlichkeit also nur in den
kleinsten Theilen Statt finden können.
Sind die Plätze, welche einer bestimmten Anzahl von gegebnen Punkten der
ersten Ebne in der Darstellung entsprechen sollen, vorgeschrieben, so kann
man leicht nach der gemeinen Interpolationsmethode die einfachste algebraische
Function / finden, wodurch diese Bedingung erfüllt wird. Bezeichnet man nem-
lich die Werthe von oc-\-iy für die gegebnen Punkte durch a,h, c n.s.w., und
die correspondirenden Werthe von X + zF durch A,B, C u. s.w. , so wird man
J " {a-b){a-c) . . .' ^ ^ {b — a){b~a).. . ' -^ ^ {c ~ a){c~b) . . .• ^ I ^ ^^•
setzen müssen, welches eine algebraische Function von u ist, deren Ordnung um
eine Einheit kleiner ist, als die Anzahl der vorgegebnen Punkte. Für zwei Punkte,
wo die Function linearisch wird, findet folglich vollkommene Aehnlichkeit Statt.
Alan kann von diesem Verfahren in der Geodäsie eine nützliche Anwendung
machen , um eine auf mittelmässige Messungen gegründete Karte, die im kleinen
Detail gut, aber im Ganzen etwas verzerrt ist, in eine bessere zu verwandeln,
wenn man die richtige Lage einer Anzahl von Punkten kennt. Es versteht sich
jedoch , dass man bei einer solchen Umformung nicht viel über die Gegend hin-
ausgehen darf, welche letztere Punkte umfassen.
Wenn man die zweite Auflösung auf dieselbe Art durchführt, so findet man,
dass der ganze Unterschied nur darin besteht , dass die Aehnlichkeit eine ver-
kehrte ist , indem alle Elemente in der Darstellung zwar eben so grosse Winkel
mit einander machen , wie im Dargestellten , aber in verkehrtem Sinn , so dass

dort rechts liegt, was hier links ist. Dieser Unterschied ist aber kein wesentli-
cher, und verschwindet , wenn man in der einen Ebne diejenige Seite, welche
man vorher als obere betrachtete , zur untern macht. Diese letztre Bemerkung
lässt sich übrigens allemal in Anwendung bringen, wenn die eine der beiden Flä-
chen eine Ebne ist , daher wir in den folgenden Beispielen dieser Art uns bloss
auf die erste Auflösung beschränken können.
9.
Wir wollen nun (als zweites Beispiel) die Darstellung der Fläche eines
geraden Kegels in der Ebne betrachten. Als Gleichung der erstem nehmen
wir an
cox-\-yy — kkzz = 0
wo wir ferner
X =. ktcosu
y r= ht^inu -
z =^t
und wie vorhin F=r, F=C7, Z=0 setzen.
Die Differentialgleichung
(1) = [kk-^\)^f-\-kktt^u' = 0
gibt hier die beiden Integrale
logf + ^V^^.w = Const.
Wir haben demnach die Auflösung
X+iF=/{log<+iv/;j^-.«)
d. i. es wird , i ndem f eine willkürliche Function bedeutet , f ür X der reelle
Theil von
/(iog*+iv/sii--")
und für F der imaginäre , nach Weglassung des Factors i, angenommen.

DIE THEILE EINER GEGEBNEN FLÄCHE AUF EINER ANDERN ABZUBILDEN ETC. 203
Setzt man für / z.B. eine Exponentialgrösse, nemlich
wo h constant ist und e die Basis d/ero h=y p ehreb' olischen Logarithmen bedeutet , so
hat man die einfachste Darstellung
Die Anwendung der Formeln des 7 . A rt. gibt hier
n = [kk^\)tt, N= 1
und, da cpu = <p u = he' ,
folglich
also constant. Macht man also noch h
h = \/{kk-}-l)
so wird die Darstellung eine vollkommne Abwicklung.
10.
Es sei drittens die Kugelfläche, deren Halbmesser = a, in der Ebne dar-
zustellen. Wir s etzen hier
a: = a cos # . s in w
j/ = asin ^.siiiM
z z=: a cos u
wodurch wir erhalten
Die Difterentialformel w = 0 gibt folglich
' sinu
und deren Integration
? + ilog cotang ^ M = Const.

Es wird daher , wenn wir wiederum durch die Charakteristik / eine will-
kürliche Funktion andeuten, X dem reellen und i Y dem imaginären Theile von
f{t -\-i log CO tang^u)
gleich gesetzt werden müssen. Wir wollen ein Paar specielle Fälle dieser allge-
meinen Auflösung anführen.
Wählt man für /" e ine lineare Function, indem man fo = A'u setzt, so wird
X=kt, Y = klogcotaug^u
Auf die Erde angewandt, ist dies, wenn man t die geographische Länge, 9ü" — u
die Breite bedeuten lässt, offenbar mit Mercators Projection einerlei. Für das
Vergrösserungsverhältniss geben hier die Formeln des 7. Artikels
k
asin«
m = —. —
Nimmt man für / eine imaginäre Exponentialfunction, und zwar zuerst die
einfachste /u = ke^"^, so wird
f{t-\-ilogcotang^u) = ke^^ ang,,u + t __ ^^angf w(cos ^-|-^■sin^)
und ^
X = A:tang4-w.cos^, F ^ Atangf w.sin^
welches, wie man leicht sieht, die stereographische Polarprojection ist.
Setzt man allgemeiner /u = ke^^"'^, so wird
X = ktang-^u'' ,cosXt, Y = kterngj^u^ .sinXt
Für das Vergrösserungsverhältniss erhalten wir hier
n ^ aasiiiu', N = \, cpu = i'kke^^-*
und hieraus
m == a smA - u —
Man sieht, dass hier die Darstellung aller Punkte, für welche u constant
ist , in Einen Kreis, und die Darstellung aller Punkte, für welche t constant ist,
in Eine gerade Linie fallt , wie auch , d ass die allen verschiedenen Werthen von
u angehörigen Kreise concentrisch sind. Dies gibt eine sehr zweckmässige Kar-
tenprojection, wenn nur ein Theil der Kugelfläche darzustellen ist, und man thut
dann am besten , X so zu wählen , dass das Vergrösserungsverhältniss für die

DIE THEILE EINER GEGEBNEN FLÄCHE AUF EINBR ANDERN ABZUBILDEN ETC. 205
äusse
^'
r sten Werthe von u gleich gross wird, wodurch es gegen die Mitte zu seinen
kleinsten Werth erhält. Sind diese äussersten Werthe von u diese m" und u,
so wird man demnach setzen müssen :
^ log sin u'— log sin u'*
logtangA«'— logtangAM»
Die Blätter von Herrn Professor Harding's Sternkarten Nr. 19 — 26 sind nach
dieser Projection gezeichnet.
11.
Man kann die allgemeine Auflösung für das im vorhergehenden Artikel be-
handelte Beispiel noch in einer andern Form aufstellen, die wir ihrer Eleganz we-
gen hier noch beifügen zu müssen glauben.
In Folge des im 6. Art. Vorgetragenen wird, da
taiig ^u {cos t-{-i sin t)
eine Function von
^ -[- ^ l og cotang 4- M
ist, und
, ^ I • • .\ sin M cos ^ -(- i sin ?< sin ^ x + iy
tang4-w(cos^-f-2smn = , _ ,' , = — r-;^
die allgemeine Auflösung auch durch
' '' a-\-z -f a + z
dargestellt werden können, d. i. X muss dem reellen und iY dem imaginären
Theil von y^+i^ gleich gesetzt werden, indem / eine willkürliche Function be-
zeichnet. Anstatt /*^4-- kann man, wie man leicht sieht, auch eine will-
kürliche Function von 0a:i5 + xoder von ^a—+ -y^ nehmen.
12.
Wir wollen viertens die Darstellung der Oberfläche des Kevolutions - E Uip-
soids in der Ebne betrachten. Es seien a und b die beiden halben Hauptaxen
des EUipsoids , so dass
0? = a cos t s in u
1/ = asiiit sin u
z =. bcosu

gesetzt werden kann. Hier wird also
u) = aasinu^ df -{- {aacosu^ -{- bb sinu^) du^
und die Differentialformel m = 0 gibt, wenn wir Kürze halber \/(l ——)= e
setzen (insofern die Revolutionshalbaxe b<^a),
0 = d#ipidw.\/(cotangw*4-l — ^^)
Setzt man hier
^(1 — gg) .tangw = tangw
wo, bei der Anwendung auf das Erdsphäroid, 90^ — w die geographische Breite
und t die Länge vorstellen wird, so verwandelt sich diese Gleichung in
deren Integration ,
Const. = ^ + üog. jcotangf2^.(|~^^3)^'}
f^ibt. Man hat daher , i ndem / eine willkürliche Function bedeutet, für X den
reellen und für iY den imaginären Theil von
/(.+ aog|eotangf...(l^:-:f!)
ZU nehmen. — Wählt man für / eine lineare Function, d. i. /u = ku. so wird
_._ ■ Tr- ^ 1 , .71 1 + ecosw
X = A-f, Y= A-logcotangJ-i^ — ^Arglog^^:::^-
welches eine der MERCATORSchen analoge Projection gibt.
Nimmt man hingegen für / eine imaginäre Exponentialfunction /u = ke^^'^,
SO wird
X==A.tang4-.^'-.(^^-)= .cosX^ F= Atang-^.'-. (^^-^ .smX^
welches, wenn man X = 1 setzt, eine der stereographischen Polarprojection ana-
loge ,u nd allgemein , eine zur Darstellung eines Theils der Erdoberfläche, inso-
fern man auf die Abplattung Rücksicht nehmen soll, sehr zweckmässige Pro-
jection gibt.
Was über den andern Fall , wo b^a ist , zu sagen ist , l ässt sich zwar
leicht aus dem vorhergehenden unmittelbar ableiten , wo , wenn man dieselben

DIE THEILE EINER GEGEBNEN FLÄCHE AUF EINER ANDERN ABZUBILDEN ETC. 207
Bezeichnungen beibehält, g imaginär, aber (7117"-—)'^ doch wieder reell wird.
Der Vollständigkeit wegen wollen wir jedoch die Formeln für diesen Fall noch
besonders beifügen, und gleich Anfangs \/(— — l) = t] setzen. Man hat dann
w durch die Gleichung
\/( 1 -|- T] T ]) . t ang w ;= tang «<?
zu bestimmen, und die Differentialgleichung
0 = d^+fd?^.— — i:tl!L_^
(1 + '^ 7 ) c os zf; ) s in 1/7
wird das Integral
Const. = t-^i (log cotang ^WJ -|-''1 ^^^ tang. tj c os w]
geben , so dass X für den reellen und ^ Y für den imaginären Theil von
/(#-j- « (log cotang^ WJ -}- T] A re tang. T] c os w) )
wird genommen werden müssen. Die Gegenstücke der beiden obigen speciellen
Anwendungen ergeben sich hieraus von selbst. Nach der erstem wird
X = kt, F= Älogcotang^-w-j-TjÄArctang.Tjcos^i'
nach der andern
X == A:tangii*;\^-^^^^^*""^-^''°^".cosX^
F = >;; t ang-i-i/ . . -''^'^^'^*^"^'^'^''^" . s in X?
gesetzt werden müssen.
13.
Als letztes Beispiel wollen wir die allgemeine Darstellung der Oberfläche
des Umdrehungs-Ellipsoids auf der Kugelfläche betrachten. Für jenes wollen wir
die Bezeichnungen des vorhergehenden Artikels beibehalten, den Halbmesser der
Kugelfläche = A, und
X== ^cosTsinL^
F= A^mTüwU , ,
Z = AcobU

setzen. Wenn man hier die allgemeine Auflösung des 5. Artikels i^ur Anwen-
dung bringt, so findet man, dass, indem / eine willkürliche Function bedeutet,
T dem reellen und Hogcotang^t^ dem imaginären Theile von
/(^+zlogjcotang-i-w).(j^^)*'i)
gleich gesetzt werden muss *). - -
Die einfachste Auflösung wird sein , yo = u zu setzen , wodurch
r=., tangil7=tang+«..(|±i£S)^«
wird. Dies bietet eine für die höhere Geodäsie überaus brauchbare Transformation
dar, von welcher Benutzung wir jedoch hier nur einiges und nur kurz andeuten
können. Wenn nemlich auf der Oberfläche des Ellipsoids und der Kugel dieje-
nigen Punkte als einander correspondirend angesehen werden, die einerlei Länge
haben, und deren Breiten resp. 90" — w, 90*^ — U, vermöge der angeführten Glei-
chung zusammenhangen, so entspricht einem System von, verhältnissmässig, klei-
nen Dreiecken (und das werden diejenigen immer sein, die zur wirklichen Mes-
sung dienen können), die auf der Oberfläche des Sphäroids durch kürzeste Linien
gebildet werden , auf der Kugelfiäche ein System von Dreiecken , deren Winkel
den correspondirenden auf dem Sphäroid genau gleich sind, und deren Seiten von
grössten Kreisboo^en so wenia: abweichen, dass sie in den meisten Fällen, wo nicht
die alleräusserste Schärfe verlangt wird , als damit zusammenfallend betrachtet
werden können , so wie auch da , wo die grösste Genauigkeit gefordert wird , d ie
Abweichung vom grössten Kreise leicht mit aller nöthigen Schärfe durch einfache
Formeln sich berechnen lässt. Man kann daher das ganze System, nachdem man
zuerst eine Dreiecksseite auf die Kugelfläche gehörig übertragen hat, ganz so, als
wenn es auf dieser selbst läge , vermittelst der Winkel berechnen , nöthigenfalls
mit der eben angedeuteten Modification, für alle Punkte des Systems die Werthe
von T und U bestimmen , und von letztern auf die correspondirenden Werthe
von w (am einfachsten vermittelst einer äusserst leicht zu construirenden Hülfs-
tafel) zurückgehen.
*) Wir übergehen hier theils die zweite Auflösung des 5. Artikels, die sich von der obigen nur
durch Vertauschung von — 2' gegen + 2' unterscheiden und einer verkehrten Darstellung entsprechen
würde, theils den Fall eines länglichen Ellipsoids, dessen Behandlung nach dem, was im vorigen Art.
vorgekommen, sich aus der des abgeplatteten von selbst ergibt.

DIE THEILE EINER GEGEBNEN FLÄCHE AUF EINER ANDERN ABZUBILDEN ETC. 209
Insofern ein Dreiecksnetz sich doch immer nur über einen sehr massigen
Theil der Erdoberfläche erstreckt, lässt sich der erwähnte Zweck noch vollkomm-
ner erreichen , wenn man die allgemeine Auflösung noch etwas generalisirt , und
nicht y*ü = ü , sondern /o = u -j- Const. annimmt. Off'enbar würde hiedurch
gar nichts gewonnen , wenn man dieser Constante einen reellen Werth beilegte,
weil dadurch lediglich T und t um diese Constante verschieden , also nur die
Anfangspunkte der Längen ungleich werden würden. Allein ganz anders verhält
es sich, wenn man der Constante einen imaginären Werth beilegt. Setzt man
dieselbe = — ilo^k, so wird
r=*, tangfr7=itang^«,.(i±i^f
Um hier über den zweckmässigsten Werth von k entscheiden zu können,
müssen wir vor allen Dingen das Vergrösserungsverhältniss bestimmen.
Es wird hier, in den Zeichen des 5. und 7. Artikels
cpu = 1
Also
m = — ;— = — ,— .vfl — ggcosir) = -.
asinw asm«) VV c, v-o iv y ^ . cos^tc*(l— e Ä c ( o l s t < — j )* e e + c o Ä 8 X M ; ? s * m ) J ^ ^ + M ^ ) *(l +ecosM))*
welches Verhältniss also bloss von der Breite abhängt. Die möglich geringste
Abweichung von voUkommner Aehnlichkeit erhält man, wenn man k so bestimmt,
dass m für die äussersten Breiten gleich grosse Werthe erhält, wodurch von selbst
m bei der mittlem Breite seinem grössten oder kleinsten Werthe sehr nahe sein
wird. Bezeichnet man die äussersten Werthe von w durch w^ und w\ so erhält
man auf diese Weise
, c os^-w)" (l — £ cos lg")' cos^w'^{\ — ecosw')*
k = i V / s i ( 1 n ^ — M > £ '' £ ' C ( O l S + t £ g CO ' S ' M V ) + ' i ) ^ ' s ( i l n ^ — M ? e » e c o ( s i t 4 g - ' £ c o ) s T m + ;» i ) ^ «
(l-££cosM;'')*+i^ (l-eecos«;<»')i+i^
Um zu erfahren , bei welcher Breite m seinen grössten oder kleinsten Werth er-
hält, haben wir

dm j TT ^ i. j I eecosw .sinw.dw
— = cot^ngU.du—cotangw.dw-{ ,_,,,^,^^
d U äw eesinw .dw {\— ee)dw
sin U sinw l — eecosw* (l — eecosto*)sinM7
und hieraus
dm ^ (l-es)d^. ^ , ^ jj _ ^^^ ^j
m sinw(l — eecost/7*) ^ '
Hieraus erhellt, dass m da seinen grössten oder kleinsten Werth erhält, wo
XJ = w wird ; b ezeichnet man den Werth von w an dieser Stelle durch W,
so wird
,1-scos^^ie ^^^^ c0sTr:==i::AL
woraus man W bestimmen kann, wenn k nach der obigen Formel berechnet
ist. Für die Ausübung wird inzwischen auf die ganz genaue Gleichheit der
Werthe von m an den äussersten Breiten wenig ankommen, und man kann sich
begnügen, für QO'' — W ungefähr die mittlere Breite zu wählen, und daraus k
abzuleiten. Den allgemeinen Zusammenhang zwischen U und w gibt dann die
Formel
t,an g4 -- T TT J = z 1ta.n g4 -, M ; (()l— — e CO S 7F ) ' ()1/4- £ COS «<-(')) 2^
O •* c3 ^ <(l + £COS JF)(l SCO&W)^
Zur wirklichen numerischen Berechnung ist es jedoch vortheilhafter , Reihen an-
zuwenden, denen man verschiedene Formen geben kann, bei deren Entwicklung
wir uns aber hier nicht aufhalten.
Da man übrigens leicht sieht, dass für w<iW, W^w, also cos C/" — cosw
und mithin auch ^ negativ; und für w^W, U<iw, mithin ^ positiv wird,
so ist klar, dass fm w := ü = W der Werth von m allemal ein Minimum wird,
und zwar
Wählt man also den Halbmesser der Kugel Ä = -Tr^zriT^^ppp^^ ^^ i^^ ^i® Dar-
stellung unendlich kleiner Theile des Ellipsoids bei der Breite 90^ — W dem Ur-
bilde nicht bloss ähnlich, sondern gleich, bei andern Breiten aber grösser.
Man kann den Logarithmen von m mit Vortheil in eine nach den Poten-
zen von cos U — cos W fortlaufende Reihe entwickeln , d eren erste für die Aus-
übung zureichende Glieder diese sind

DIE THEILE EINER GEGEBNEN FLÄCHE AUF EINER ANDERN ABZUBILDEN ETC. 211
loghyp.m =. log {^V(l -escos W^)\-\-~^^ . ( cos U— cos Wf
_i!:^.(cos?7-coslF)\...
3(1— ee)* ^ ''
Wenn also z. B . die Dänische Monarchie innerhalb der Grenzen der Breite
53" und 58*^ auf diese Weise auf die Kugelfläche übertragen und W = 34*^30'
gesetzt wird, so wird bei der Abplattung -3-^^ die Darstellung an den Grenzen,
linearisch gerechnet , nur um ^ 3 ^ ^ 0 0 vergrössert.
Wir müssen uns hier damit begnügen , nur eine kurze Andeutung von ei-
ner Benutzungsart des Uebertragens der Figuren in der höhern Geodäsie gegeben
zu haben , und eine angemessenere Ausführung für einen andern Ort versparen.
14.
Es bleibt uns noch übrig , einen in unsrer allgemeinen Auflösung vorkom-
menden Umstand hier etwas ausführlicher zu betrachten. Wir haben im 5. Ar-
tikel gezeigt, dass allemal zwei Auflösungen statt finden, indem entweder P-f-iQ
einer Function von p-\-iq, und P — iQ einer Function von p — iq gleich wer-
den muss; oder P-\-iQ einer Function von p — iq, und P — iQ einer Function
von p-\-iq. Wir wollen nun noch zeigen, dass allemal bei der einen Auflösung
die Theile in der Darstellung zugleich eine ähnliche Lage haben , wie im Darge-
stelltenb ; ei der andern Auflösung hingegen verkehrt liegen ; zugleich wollen wir
dasCriterium angeben, nach welchem dieses a priori unterschieden werden kann.
Zuvörderst bemerken wir, dass von vollkommner oder verkehrter Aehnlich-
keit nur insofern die Rede sein kann, als an jeder der beiden Flächen zwei Seiten
unterschieden werden, wovon die eine als die obere, die andere als die untere
betrachtet wird. Da dieses an sich etwas willkürliches ist, so sind beide Auflö-
sungen gar nicht wesentlich verschieden , und eine verkehrte Aehnlichkeit wird
zur vollkommnen , sobald man bei der einen Fläche die vorher als obere betrach-
tete Seite zur untern macht. Bei unsrer Auflösung konnte daher diese Unterschei-
dung gar nicht vorkommen, da die Flächen bloss durch die Coordinaten ihrer
Punkte bestimmt wurden. Will man auf diesen Unterschied eingehen, so muss
zuvor die Natur der Flächen auf eine andere Art festgelegt werden , welche ihn
mit in sich fasst. Zu diesem Zweck wollen wir annehmen, dass die Natur der er-
sten Fläche durch die Gleichung cj; = 0 bestimmt werde , wo ^ eine gegebne
einförmige Function von x, y, z ist. In allen Punkten der Fläche wird also der
35*

ALLGEMEINE AUFLÖSUNG DER AUFGABE
Werth von d> verschwinden, und in allen Punkten des Raumes, welche der Fläche
nicht angehören, wird er nicht verschwinden. Bei einem Durchgange durch die
Fläche wird also , wenigstens allgemein zu reden, der Werth von (|> a us dem Po-
sitiven ins N egative , b ei dem entgegengesetzten aus dem Negativen ins Positive
übergehen , oder auf der einen Seite der Fläche wird der Werth von cf» p ositiv,
auf der andern negativ sein : die erstere wollen wir als die obere , die andere als
die untere betrachten. Ganz eben so soll es bei der zweiten Fläche gehalten wer-
den ,i ndem ihre Natur durch die Gleichung W = 0 bestimmt wird, wo ¥ eine
gegebne einförmige Function der Coordinaten X, Y, Z ist. Es gebe ferner die
Differentiation
d(j; = edx -j-^dy -\-hdz
d'^=EdX-{-GdY-{-HdZ
wo e,g,h Functionen von x,y,z und JE, G,Il Functionen von X, F, Z
sein werden.
Da die Betrachtungen , durch welche wir zu dem vorgesetzten Ziele gelan-
gen müssen, obwohl an sich nicht schwierig, doch etwas ungewöhnlicher Art sind,
so wollen wir uns bemühen, ihnen die grösste Klarheit zu geben. Wir wollen
zwischen den beiden einander entsprechenden Darstellungen auf den Flächen,
deren Gleichungen (J> = 0 und »F = 0 sind , sechs Zwischen-Darstellungen in
der Ebne annehmen , so dass acht verschiedene Darstellungen in Betracht kom-
men, nemlich
P,
indem als co.r r Xe,sp o nyd,i rend be-
trachtet werden die Punkte,
deren Coordinaten resp. =
das Urbild in der Fläche , deren Gleichung c|; = 0
Darstellung in der Ebne oe,
t.
r, V, 0
X, F, 0
Abbildung in der Fläche, deren Gleichung ¥ = 0 X, F,Z
Wir wollen nun diese verschiednen Darstellungen unter einander lediglich
in Beziehung auf die gegenseitige hage der unendlich kleinen Linearelemente ver-

DIE THEILE EINER GEGEBNEN FLÄCHE AUF EINER ANDERN ABZUBILDEN ETC. 213
gleichen, indem wir das Grössenverhältniss ganz bei Seite setzen; als ähnlichlie-
gend werden also zwei Darstellungen betrachtet, wenn von zwei aus Einem
Punkte ausgehenden Linearelementen dem in der einen Darstellung rechts liegen-
den auch in der andern das rechts liegende entspricht; im entgegengesetzten Falle
werden sie verkehrtliegende heissen. Bei der Ebne, von Nro. 2 — 7 wird immer
die Seite , wo die positiven Werthe der dritten Coordinate liegen , als die obere
betrachtet; bei der ersten und letzten Fläche hingegen ist die Unterscheidung der
obern und untern Seite bloss von dem positiven oder negativen Werthe von
^ und ¥ abhängig , wie schon oben festgesetzt ist.
Hier ist nun zuvörderst klar, dass für jede Stelle der ersten Fläche, wo
man bei ungeändertem x und i/ durch ein positives Increment von z auf deren
obere Seite kommt, die Darstellung in 2 mit der in 1 ähnlichliegend sein wird ;
dies wird also offenbar überall zutreffen, wo h positiv ist; und das Gegentheil
wird bei einem negativen h eintreten, wo die Darstellungen verkehrt liegend sein
werden.
Auf dieselbe Weise werden die Darstellungen in 7 und 8 ähnlich liegend
oder verkehrt liegend sein , j enachdem H positiv oder negativ ist.
Um die Darstellungen in 2 und 3 unter sich zu vergleichen , sei in der er-
stem d^ die Länge einer unendlich kleinen Linie von dem Punkte, dessen Coor-
dinaten ^, y, zu einem andern, dessen Coordinaten ic-\-da!, i/-\-di/ sind, und
/ dessen Neigung gegen die Abscissenlinie wachsend in dem Sinn , i n welchem
man von der Axe der cc zu der Axe der ^ übergeht, also
dod = ds .cosl, d^ = ds.sml
In der Darstellung 3 sei do die Grösse der Linie, welche der ds entspricht, und
ihre Neigung zur Abscissenlinie , wie vorhin verstanden , X , so dass
d^ = da.cosX, dw = do.sinX
Man hat also, in den Bezeichnungen des 4. Artikels
ds.cosZ= da(acosX-|-a'sinX)
dÄ . s in ^ = d 0 ( 6 c os X -|- 6'sinX)
folglich
tang/ == '^"\^t':^!\'

Betrachtet man nun x und y als constant, und /, X als veränderlich, so gibt
die Differentiation
d^ ah' — ha / ii l.„'\ ^_^
dX (oco8X + a'8inX)*+(öcosX + 6'8inXf ^""^ ^^^ds*
Man sieht also, dass jenachdem ah' — ha positiv 06?er negativ ist , / und X im-
mer zugleich wachsen, oder sich entgegengesetzt ändern, und also im erstem
Fall die Darstellungen 2 und 3 ähnlich liegend, im andern verkehrt liegend sind.
Aus der Verbindung dieses Resultats mit dem vorhingefundenen ergibt sich,
dass die Darstellungen in 1 und 3 ähnlich liegend oder verkehrt liegend sind , j e
nachdem "^"^"^ positiv oder negativ ist.
Da auf der Fläche, deren Gleichung c); = 0 ist,
ediOO-\-gdy-{-hdiZ =^ 0
also auch
[ea-\-gh-\-hc)6.t-\-{ed-\-gb'-\-hc)diU =^ 0
wird, wie auch immer das Verhältniss von dt und dw gewählt wird, so muss
offenbar identisch
ea-{-gh-\-hc = 0, ed-\-gV-\-hc =^ 0
werden, woraus folgt, dass e, g, ä resp. den Grössen hc — ch\ cd — ac\ ah' — ha
proportional sind , also
bc' — ch' cd — ac' ah' — ha'
e g Ä
Man kann also , welchen dieser drei Ausdrücke man will , oder wenn man mit
der ihrer Natur nach positiven Grösse ee-\-gg-\-hh multiplicirt , d ie sich erge-
bende symmetrische Grösse
ehc'-\-gcd-\-hah' — ech' — gac — hhd
als Criterium der ähnlichen oder verkehrten Lage der Theile in den Darstellun-
gen 1 u nd 3 anwenden.
Ganz eben so wird ähnliche oder verkehrte Lage der Theile in den Darstel-
lungen 6 u nd 8 von dem positiven oder negativen Werthe der Grösse
BC'-CB' _ CA'-AC AB'-£A'
E ~ G ~~' H

DIE THEILE EINER GEGEBNEN FLÄCHE AUF EINER ANDERN ABZUBILDEN ETC. 215
oder 'wenn man lieber will , der symmetrischen
EBC'-{-GCÄ-\-HÄB'—ECB'—GÄC'—HBÄ'
abhangen.
Die Vergleichung der Darstellungen in 3 und 4 beruhet auf ganz ähnlichen
Gründen , wie die von 2 und 3 , und die ähnliche oder verkehrte Lage der Theile
hängt von dem positiven oder negativen Zeichen der Grösse
ab ; und eben so bestimmt das positive oder negative Zeichen von
^dTJ'^dUJ ^dU^^dT'
die ähnliche oder verkehrte Lage der Theile in den Darstellungen 5 und 6.
Was endlich die Vergleichung der Darstellungen 4 und 5 unter sich be-
trifft, so k önnen wir uns auf die Analyse des 8. Artikels beziehen, aus welcher
erhellet, dass jene in den kleinsten Theilen ähnlich, oder verkehrt liegend sind,
je nachdem man die erste oder zweite Auflösung gewählt , d . i. entweder
P-{-iQ=f[p-^iq) und P—iQ=f{p — iq)
oder
P-\-iQ=f[p — iq) und P—iQ=f{p-{-iq)
gesetzt hat.
Aus diesem allen ziehen wir nunmehro den Schluss , d ass man , wenn die
Darstellung auf der Fläche , deren Gleichung W = o ist , dem Urbilde auf der
Fläche , deren Gleichung c|; = 0 ist , i n den kleinsten Theilen nicht bloss ähn-
lich, sondern auch ähnlich liegend sein soll, auf die Anzahl der negativen Grössen,
welche unter diesen vier Grössen vorkommen ,
Ä ' ^ dt''^dil~^d^''\Jil' ^dT/'^dÜI ^dTJ''^dTl' H
Rücksicht nehmen muss ; i st gar keine oder eine gerade Anzahl darunter, so wird
die erste ; i st eine oder drei negative unter ihnen , so wird die zweite Auflösung
gewählt werden müssen. Bei entgegengesetzter Wahl findet allemal eine ver-
kehrte Aehnlichkeit Statt.

Uebrigens lässt sich noch zeigen , dass , wenn obige vier Grössen resp. mit
r, s, S, R bezeichnet werden , allemal
rs/{ee + gg + hh) ^^^ RUEE+GG^HH^ ^^^
s — ' ^ —
wird, n und N in der Bedeutung des 5. Art. genommen; wir übergehen jedoch
hier den nicht schwer zu findenden Beweis dieses Theorems , da dieses für un-
sern Zweck nicht weiter nöthig ist. . '•
[Randbemerkungen in Gkvs,^ Handexemplar ■]
[Art. 10 neben der letzten Gleichung zur Bestimmung von X] oder X = cosm*, wenn für u = u*
der Minimalwerth [des Vergrösserungsverhältnisses] Statt finden soll.
[Art. 1 2 neben der Gleichung, durch welche hier im Abdruck die Grösse w eingeführt wird] Das Zei-
chen 0) i st gegen meine Absicht im Druck gebraucht : es sollte w sein.
[Art. 1 3 neben den Gleichungen , die sich auf die durch die Function /u = u — «log k bestimmte Ab-
bildung beziehen, sind die entsprechenden Gleichungen für die Function /u = au — tlogÄ verzeichnet,
welche später in der ersten Abhandlung der Untersuchungen über Gegenstände der höhern Geodäsie aufge-
nommen wurden.]


\newpage
\selectlanguage{latin}
\section*{Disquisitiones generales circa superficies curvas}
\addcontentsline{toc}{section}{Disquisitiones generales circa superficies curvas}

\begin{center}
\textsc{Auctore}\\
CAROLO FRIDERICO GAUSS\\
\smallskip
SOCIETATI REGIAE OBLATAE D.\ VIII.\ OCTOBR.\ MDCCCXXVII.\\
\smallskip
{\small Commentationes societatis regiae scientiarum Gottingensis recentiores. Vol.~vi.\\
Gottingae mdcccxxviii.}
\end{center}

DISQUISITIONES GENERALES
CIRCA SUPERFICIES CURVAS
A U C T 0 R E
CAROLO FEIDEEICO GAUSS
SOCIETATI REGIAE OBLATAE D. Vlll. OCTOBR. MDCCCXXVIl.
Commentationes societatis regiae scientiarum Gottingensis recentiores. Vol. vi.
Gottingae mdcccxxviii.



DISQUISITIONES GENERALES
CIRCA SUPERFICIES CURVAS.
1.
Disquisitiones , i n quibus de directionibus variarum rectarum in spatio agi-
tur , plerumque ad maius perspicuitatis et simplicitatis fastigium evehuntur , i n
auxilium vocando superficiem sphaericam radio = 1 circa centrum arbitrarium
descriptam, cuius singula puncta repraesentare censebuntur directiones rectarum
radiis ad illa terminatis parallelarum. Dum situs omnium punctorum in spatio
per tres coordinatas determinatur, puta per distantias a tribus planis fixis inter
se normalibus, ante omnia considerandae veniunt directiones axium bis planis
normalium : puncta superiiciei sphaericae, quae has directiones repraesentant, per
(l), (2), (3) denotabimus; mutua igitur horum distantia erit quadrans. Ceterum
axium directiones versus eas partes acceptas supponemus, versus quas coordina-
tae respondentes crescunt.
2.
Haud inutile erit, quasdam propositiones , quae in huiusmodi quaestioni-
bus usum frequentem ofFerunt, hie in conspectum producere.
I. Angulus inter duas rectas se secantes mensuratur per arcum inter puncta,
quae in superficie sphaerica illarum directionibus respondent.
II. Situs cuiuslibet plani repraesentari potest per circulum maximum in
superficie sphaerica, cuius planum illi est parallelum.
36=^

III. Angulus inter diio plana aequalis est angulo sphaerico inter circulos
maximos illa repraesentantes , et proin etiam per arcum inter herum circulorum
maximorum polos interceptum mensuratur. Et perinde inclinatio rectae ad pla-
num mensuratur per arcum , a puncto , quod respondet directioni rectae , ad cir-
culum maximum, qui plani situm repraesentat , normaliter ductum.
IV. Denotantibus x,y,z; x\ y\ z coordinatas duorum punctorum , r eo-
rundem distantiam, atque L punctum, quod in superficie sphaerica repraesen-
tat directionem rectae a puncto priore ad posterius ductae , erit
a?' = a? + rcos(l)X, y = y -{-rcos[1)L, 2;' = 0-|-rcos(3)i
V. Hinc facile sequitur, haberi generaliter
cos(l)i.'4-cos(2)i.'4-cos(3)X' = 1
nee non , d enotante L' quodcunque aliud punctum superficiei sphaericae , esse
cos (1 ) i . c os (1 ) L '-\- cos (2) L . c os (2) L -{- cos (3) L . c os {Z)L' = cos LL'
VI. Theorema. Denotantibus L, L', L", L'" quatuor puncta in superficie
sphaerae, atque A angulum, quem arcus LH, L'L" in puncto concursus sui for-
mant, erit
cos LL". cos L'L'"— cos LL'". cos L'L" = sm LL'. sin L"L'". cosA
Dem. Denotet litera A insuper punctum concursus ipsum , statuaturque
AL = t, AL' = t', AL" = t", AL'"= t'"
Habemus itaque :
cosLL" = costcost" -\-smtsmt"cosA
cos L'L" = cos ^'cos t'"-{- sin t'sm t'"cos A
cos L L" = cos t c os t'"-^ sin t s in t'"cos A
cosL'L" = cos?'cos#"-l-sin^'sin^"cos J.
et proin
cos LL". cos LL" — cosLL'". cos L'L" '
:= COS A (cos t c os t" siu t' s in t'"-\- cos t'cos t'" s in t s in t"
— cos ^cos Tsin t'sint" — cos ^'cos ^"sin ?sin f)
= cos A (cos t s in t'— sin t c os t' ) ( cos ^"sin t'" — sin t " c os t'" )
.= cos^.sin(^' — t).sm{t'" — t")
z= cos A. sin LL'. sin L'L"

CIRCA SUPERFICIES CURVAS.
Ceterum quum inde a puncto A bini rami utriusque circuli maximi proficiscan-
tur , duo quidem ibi anguli formantur , quorum alter alterius complementura ad
180": sed analysis nostra monstrat, eos ramos adoptandos esse, quorum directio-
nes cum sensu progressionis a puncto L ad L\ et a puncto L" ad L"' consen-
tiunt: quibus intellectis simul patet, quum circuli maximi duobus punctis con-
currant, arbitrarium esse, utrum eligatur. Loco anguli Ä etiam arcus inter po-
los circulorum maximorum, quorum partes sunt arcus LL\ L"L"\ adhiberi pot-
est: manifeste autem polos tales accipere oportet, qui respectu herum arcuum si-
militer iacent, puta vel uterque polus ad dextram iacens, dum a L versus L'
atque ab L" versus L" procedimus , vel uterque ad laevam.
VII. Sint L, L\ L" tria puncta in superficie sphaerica, statuamusque bre-
vitatis caussa
cos(l)i ==07, cos(2)i =y, cos(3)Z =z
cos{i)L' = üs', cos(2)i' =:/, cos{Z)L' = z'
cos (1 ) L " = x", cos (2) L " = y" , cos (3) U ' = z"
nee non
xy'z'-\- xy"z -\- oc'y z — xy"z — xy z" — x'y'z = A
Designet X polum circuli maximi , cuius pars est arcus LL', et quidem eum, qui
respectu huius arcus similiter iacet, ac punctum (1) respectu arcus (2) ( 3). Tunc
erit, ex theoremate praecedente, yz — y^; = cos(l)X.sin(2)(3).sinXi^', sive,
propter (2) (3) = 90",
yz — y'z =^ cos(l)X.sinZyiv', et perinde
zx — z'x = cos(2)X.sinZ/X'
xy — x'y = cos(3)X.sinZ/iy'
Multiplicando has aequationes resp. per x", y", z" et addendo , obtinemus adiu-
mento theorematis secundi in Y prolati
A = cos'kL".smLL'
lam tres casus sunt distinguendi. Primo, quoties L" iacet in eodem circulo
maximo, cuius pars est arcus LL', erit XX" = 90", adeoque A = 0. Quoties
vero L" iacet extra circulum illum maximum, aderit casus seciindus, si est ab ea-
dem parte , a qua est X, tertivs , si ab opposita : in bis casibus puncta L, L\ L"

formabunt triangulum sphaericum, et quidem iacebunt in casu secundo eodem
ordine quo puncta (l), (2), (3), in casu tertio vero ordine opposito. Denotando
angulos illius trianguli simpliciter per L, L\ L", atque perpendiculum in super-
ficie sphaerica a puncto L' ad latus LL' ductum per p, erit
sinj? = miL.smLL" =^ sinL'.smL'L", atque XX"= 90^'^p
valente signo superiori pro casu secundo, inferiori pro tertio. Hinc itaque colligimus
+ A = sin i . s iniv-L'. sin -LZ" = sinL'. sin LL\ sin L'L" = sinL". sin LL". sin L'L"
Ceterum manifesto casus primus in secundo vel tertio comprehendi censeri potest,
nulloque negotio perspicitur , + A exhibere sextuplum soliditatis pyramidis in-
ter puncta L, L\ L" atque centrum sphaerae formatae. Denique hinc facillime
coUigitur, eandem expressionem +^A generaliter exprimere soliditatem cuius-
vis pyramidis inter initium coordinatarum atque puncta quorum coordinatae sunt
X, 1/, z; <^'', y , z; cc\ y" , z", contentae.
3.
Superficies curva apud punctum A in ipsa situm curvatura continua gaudere
dicitur , si directiones omnium rectarum ab A ad omnia puncta superficiei ab A
infinite parum distantia ductarum infinite parum ab uno eodemque piano per A
transiente deflectuntur: hoc planum superficiem curvam in puncto A tangere di-
citur. Quodsi huic conditioni in aliquo puncto satisfieri nequit, continuitas cur-
vaturae hie interrumpitur , uti e. g. evenit in cuspide coni. Disquisitiones prae-
sentes ad tales superficies curvas, vel ad tales superficiei partes, restringentur, in
quibus continuitas curvaturae nuUibi interrumpitur. Hie tantummodo observa-
mus, methodos, quae positioni plani tangentis determinandae inserviunt, pro
punctis singularibus , i n quibus continuitas curvaturae interrumpitur , vim suam
perdere , et ad indeterminata perducere debere.
Situs plani tangentis commodissime e situ rectae ipsi in puncto A normalis
cognoscitur, quae etiam ipsi superficiei curvae normalis dicitur. Directionem hu-
ius normalis per punctum L in superficie sphaerae auxiliaris repraesentabimus,
atque statuemus

CIRCA SUPERFICIES CURVAS. 223
COS(l)i^=rX, C0S(2)iy=F, COS(^)L = Z
coordinatas puncti Ä per .37, j/, z denotamus. Sint porro aj-{-dx, i/-\-di/, ^-f-d^
coordinatae alius puncti in superficie curva Ä'; ds ipsius distantia infinite parva
ab A; denique X punctum superficiei sphaericae repraesentans directionem ele-
menti ÄÄ. Erit itaque
da? = dÄ.cos(l)X, d^ = dÄ.cos(2)X, d^r = dÄ.cos(3)X
et, quum esse debeat Xi = 90^
Xcos(l)X4- Fcos(2)X+Zcos(3)X = 0
E combinatione harum aequationum derivamus ^
Xda^^Yd^-\-Zdz= 0
Duae habentur methodi generales ad exhibendam indolem superficiei cur-
vae. Methodus j?nma utitur aequatione inter coordinatas x,i/,z, quam reductam
esse supponemus ad formam W = 0 , ubi W erit functio indeterminatarum
X, y, z. Sit difFerentiale completum functionis W
dW=Pdx-\-Qdj/-i-Rdz
eritque in superficie curva
Pdx-{-Qdi/-{-Rdz = 0
et proin
Pcos(1)X-l-Qcos(2)X-f-J2cos(3)X = 0
Quum haec aequatio , perinde ut ea quam supra stabilivimus , valere debeat pro
directionibus omnium elementorum d^ in superficie curva, facile perspicie-
mus, X, F, Z proportionales esse debere ipsis P, Q, R et proin, quum fiat
XX-\-YY-\-ZZ=^l
erit vel
X=- , ^ Y- \ / ( P P + Q QQ + RRy ^r7 — ^ { P P + Q Q R + RR)
vel
X- s /{ P P + Q - Q ^ - \-ItM)' Y— > J { P P + - QQ Q + RRy - 7 ^ s l ^ P P ^ r - Q ^ Q ^RPt)

224 DISQUISITIOXES GENERALES
Methodus secunda sistit coordinatas in forma functionum duarum variabi-
lium p, q. Supponamus per diiferentiationem harum functionum prodire
da? = adip-{-adq
dj/ == hdip-\-h'diq
d.z =^ c6.p-\-cd.q
quibus valoribus in formula supra data substitutis , obtinemus
(aX+6F+cZ)djü-f (a'X+6'FH-c'Z)d^ = 0
Quum haec aequatio locum habere debeat independenter a valoribus differentia-
lium dp, dg, manifeste esse debebit
«X+6F+cZ=0, dX-{-h'Y-\-cZ =0
unde colligimus , X , Y, Z proportionales esse debere quantitatibus
hc — c6', ca — ac, ab' — ha
Statuendo itaque brevitatis caussa
s^([hd—chJ^[cd—acf-\-[ab'—ha:f) = ^
erit vel
-^ hc' — cV fT' ca' — ac' r/ aV — ha'
X = — ^-, y = —3—' ^ — — Ä"
vel
-er cV — hc' y;r ac' — ca' ry ha' — ah'
A = — 3— , y = -^— , Zj — — ^
His duabus methodis generalibus accedit tertia, ubi una coordinatarum,
e. g. z exhibetur in forma functionis reliquarum x, y : haec methodus manifeste
nihil aliud est, nisi casus specialis vel methodi primae, vel secundae. Quodsi
hie statuitur
d^; = tdx-{~udi/
erit vel
V __ —l y __ —^ y __ }
vel \/(i + tt + tiu)- ^(^i^tt^uu)' sj(\+tt+uu)
X = -rTrAi-. ., Y=^,-^^, ,, Z =
\J{\-\-tt-\-uu)' v^(i+^^ + «m)' sl{\-^tt + tiu)

CIRCA SUPERFICIES CURVAS. 225
5.
Duae solutiones in art, praec. inventae manifesto ad puncta superficiei
sphaericae opposita, sive ad directiones oppositas referuntur, quod cum rei natura
quadrat, quum normalem ad utramvis plagam superficiei curvae ducere liceat.
Quodsi duas piagas , superficiei contiguas , i nter se distinguere, alteramque exte-
riorem alteram interiorem vocare placet, etiam utrique normali suam solutionem
rite tribuere licebit adiumento theorematis in art. 2 (VII) evoluti, simulatque cri-
terium stabilitum est ad plagam alteram ab altera distinguendam.
In methodo prima tale criterium petendum erit a signo valoris quantitatis
W. Scilicet generaliter loquendo superficies curva eas spatii partes , in quibus
W valorem positivum obtinet, ab iis dirimet, in quibus valor ipsius W fit ne-
gativus. Et heoremate illo vero facile colligitur, si W valorem positivum obti-
neat versus plagam exteriorem, normalisque extrorsum ducta concipiatur, solu-
tionem priorem adoptandam esse. Ceterum in quovis casu facile diiudicabi-
tur , utrum per superficiem integram eadem regula respectu signi ipsius W va-
leat , an pro diversis partibus diversae : quamdiu coefiicientes P, Q, II valores
finitos habent , nee simul omnes tres evanescunt , l ex continuitatis vicissitudinem
vetabit.
Si methodum secundam sequimur , i n superficie curva duo systemata linea-
rum curvarum concipere possumus , alterum , pro quo p est variabilis , q con-
stans; alterum, pro quo q variabilis, p constans: situs mutuus harum linearum
respectu plagae exterioris decidere debet, utram solutionem adoptare oporteat.
Scilicet quoties tres lineae, puta ramus lineae prioris systematis a puncto A pro-
ficiscens crescente p, ramus posterioris systematis a puncto A egrediens crescente
q, atque normalis versus plagam exteriorem ducta similiter iacent , ut , inde ab
origine abscissarum, axes ipsarum x,y, z resp. (e. g . si tum e tribus lineis il-
lis, tum e tribus bis, prima sinistrorsum , secunda dextrorsum, tertia sursum di-
recta concipi potest), solutio prima adoptari debet; quoties autem situs mutuus
trium linearum oppositus est situi mutuo axium ipsarum 'x, y , z, solutio secunda
valebit.
In methodo tertia dispiciendum est, utrum, dum z incrementum positivum
accipit, manentibus x ei y invariatis, transitus fiat versus plagam exteriorem an
interiorem. In casu priore , pro normali extrorsum directa , solutio prima valet,
in posteriore secunda.

6.
Sicuti, per translatam directionem normalis in superficiem curvam ad su-
perficiem sphaerae, cuius puncto determinato prioris superficiei respondet punctum
determinatum in posteriore, ita etiam quaevis linea, vel quaevis figura in illa re-
praesentabitur per lineam vel figuram correspondentem in hac. In comparatione
duarum figurarum hoc modo sibi mutuo correspondentium , quarum altera quasi
imago alterius erit, duo momenta sunt respicienda , alterum , quatenus sola quan-
titas consideratur , a lterum , quatenus abstrahendo a relationibus quantitativis so-
lum situm contemplamur.
Momentum primum basis erit quarundam notionum , quas in doctrinam de
superficiebus curvis recipere utile videtur. Scilicet cuilibet parti superficiei cur-
vae limitibus determinatis cinctae curvaturam totalem seu integram adscribemus,
quae per aream figurae illi in superficie sphaerica respondentem exprimetur. Ab
hac curvatura integra probe distinguenda est curvatura quasi specifica, quam nos
mensuram curvaturae vocabimus : haec posterior ad punctum superficiei refertur, et
denotabit quotientem qui oritur, dum curvatura integra elementi superficialis
puncto adiacentis per aream ipsius elementi dividitur , et proin indicat rationem
arearum infinite parvarum in superficie curva et in superficie sphaerica sibi mu-
tuo respondentium. Utilitas harum innovationum per ea, quae in posterum a no-
bis explicabuntur , abunde , ut speramus , sancietur. Quod vero attinet ad ter-
minologiam , i mprimis prospiciendum esse duximus , ut omnis ambiguitas arcea-
tur , quapropter haud congruum putavimus , analogiam terminologiae in doctrina
de lineis curvis planis vulgo receptam (etsi non omnibus probatam) stricte sequi,
secundum quam mensura curvaturae simpliciter audire debuisset curvatura, cur-
vatura integra autem amplitudo. Sed quidni in verbis faciles esse liceret, dum-
modo res non sint inanes , neque dictio interpretationi erroneae obnoxia?
Situs figurae in superficie sphaerica vel similis esse potest situi figurae
respondentis in superficie curva, vel oppositus (in versus); casus prior locum ha-
bet ,u bi binae lineae in superficie curva ab eodem puncto directionibus inaequa-
libus sed non oppositis proficiscentes repraesentantur in superficie sphaerica per
lineas similiter iacentes , puta ubi imago lineae ad dextram iacentis ipsa est ad
dextram; casus posterior, ubi contrarium valet. Hos duos casus per Signum men-
surae curvaturae vel positivum vel negativum distinguemus. Sed manifesto haec
distinctio eatenus tantum locum habere potest, quatenus in utraque superficie pla-

CIRCA SUPERFICIES CURVAS. ■ 227
gam determinatam eligimus , i n qua figura concipi debet. In sphaera auxiliari
semper plagam exteriorem , a centro aversam , adhibebimus : in superficie curva
etiam plaga exterior sive quae tamquam exterior consideratur, adoptari potest, vel
potius plaga eadem , a qua normalis erecta concipitur ; manifesto enim respectu
similitudinis figurarum nihil mutatur, si in superficie curva tum figura ad plagam
oppositam transfertur, tum normalis, dummodo ipsius imago semper in eadem
plaga superficiei sphaericae depingatur.
Signum positivum vel negativum , quod pro situ figurae infinite parvae men-
surae curvaturae adscribimus, etiam ad curvaturam integram figurae finitae in su-
perficie curva extendimus. Attamen si argumentum omni generalitate amplecti
suscipimus , quaedam dilucidationes requiruntur , quas hie breviter tan tum attin-
gemus. Quam diu figura in superficie curva ita comparata est, ut singulis punctis
intra ipsam puncta diversa in superficie sphaerica respondeant, definitio ulteriore
explicatione non indiget. Quoties autem conditio ista locum non habet, necesse
erit , quasdam partes figurae in superficie sphaerica bis vel pluries in computum
ducere , unde , pro situ simili vel opposito , vel accumulatio vel destructio oriri
poterit. Simplicissimum erit in tali casu , f iguram in superficie curva in partes
tales divisam concipere , quae singulae per se spectatae conditioni illi satisfaciant,
singulis tribuere curvaturam suam integram , quantitate per aream figurae in su-
perficie sphaerica respondentis , signo per situm determinatis, ac denique figurae
toti adscribere curvaturam integram ortam per additionem curvaturarum integra-
rum, quae singulis partibus respondent. Generaliter itaque curvatura integra
figurae est =Jkda, denotante da elementum areae figurae , k mensuram cur-
vaturae in q uovis puncto. Quod vero attinet ad repraesentationem geometricam
huius integralis, praecipua huius rei momenta ad sequentia redeunt. Peripheriae
figurae in superficie curva (sub restrictione art. 3) s emper respondebit in superfi-
cie sphaerica linea in se ipsam rediens. Quae si se ipsam nullibi intersecat, to-
tam superficiem sphaericam in duas partes dirimet, quarum altera respondebit
figurae in superficie curva , et cuius area , positive vel negative accipienda, prout
respectu peripheriae suae similiter iacet ut figura in superficie curva respectu suae,
vel inverse , exhibebit posterioris curvaturam integram. Quoties vero linea ista
se ipsam semel vel pluries secat , exhibebit figuram complicatam , cui tamen area
certa aeque legitime tribui potest , ac figuris absque nodis, haecque area, rite in-
tellecta , semper valorem iustum curvaturae integrae exhibebit. Attamen uberio-
37*

rem huius argumenti de figuris generalissime conceptis expositionem ad aliam oc-
casionem nobis reservare debemus.
7.
Investigemus iam formulam ad exprimendara mensuram curvaturae pro quo-
vis puncto superficiei curvae. Denotante da aream elementi huius superficiei,
Zd a erit area proiectionis huius elementi in planum coordinatarum o?, y ; et per-
inde, si d2 est area elementi respondentis in superficie sphaerica, erit Zdl area
proiectionis ad idem planum : signum positivum vel negativum ipsius Z vero in-
dicabit situm proiectionis similem vel oppositum situi elementi proiecti: manifesto
itaque illae proiectiones eandem rationem quoad quantitatem , simulque eandem
relationem quoad situm, inter se tenent, ut elementa ipsa. Consideremus iam
elementum trianguläre in superficie curva, supponamusque coordinatas trium
punctorum , quae formant ipsius proiectionem , esse
0), y
xA^hx, y-\-^y
Duplex area huius trianguli exprimetur per formulam
dx.hy — dy .hx
et quidem in forma positiva vel negativa , prout situs lateris a puncto primo ad
tertium respectu lateris a puncto primo ad secundum similis vel oppositus est si-
tui axis coordinatarum y respectu axis coordinatarum x.
Perinde si coordinatae trium punctorum , quae formant proiectionem ele-
menti respondentis in superficie sphaerica , a centro sphaerae inchoatae , sunt
X, Y
X-i-dX, F+dF
X-}-gX, Y-\-lY
duplex area huius proiectionis exprimetur per
dX.^F-dF.aX
de cuius expressionis signo eadem valent quae supra. Quocirca mensura curva-

CIRCA SUPERFICIES CURVAS. 229
turae in hoc loco superiiciei ciirvae erit
, dX.oF-dF.oX
dx .hy — dy .Sa;
Quodsi iam supponimus, indolem superficiei curvae datam esse secundum modiim
tertium in art. 4 consideratum, habebuntur X et F in forma functionum quan-
titatum <r, 1/ , unde erit
d«Fr==((|^f))d«^^ + (+ i J()^8yf )dy ■■ ■
Substitutis bis valoribus , expressio praecedens transit in hanc :
, _ /dX..dF. /dX.,dF
Statuendo ut supra
ds , dz
dx ' dy
atque insuper
ddz rp ddz j^j ddz j-r
sive dx^ ^' dar.dy — ^' d^ — ^
df = rd<2?+ Ud^, du = Uda^-\- Fdy
habemus ex formulis supra datis
X = —tZ, Y=.^uZ, {l-}-tt-{-uu)ZZ= l
atque hinc
dX = —Zdt — tdZ
dY= —Zdu — udZ
{l-{-tt-^uu)dZ-{-Z{tdt-\-udu) = 0
sive
dZ = —Z\tdt-\-udu)
dX= ~Z'{l-]-uu)dt-^Z'tudu
dF = ^Z^tndt — Z'{l-^tt)du

adeoque
^^ = z'(-n+uu]u+tur)
^I = . Z'{t»T-H+tt)U)
^ = Z^tnU-{l + tt)V)
quibus valoribus in expressione praecedente substitutis , prodit
TV—UU
k = Z'{TV— ÜU){i-i-tt-^uu) = Z'iTV— Uü) = {\ + tt-\-uuf
Per idoneam electionem initii et axium coordinatarum facile effici potest, ut
pro puncto determinato A valores quantitatum t, u , U evanescant. Scilicet duae
priores conditiones iam adimplentur , si planum tangens in hoc puncto pro piano
coordinatarum ,r, 1/ adoptatur. Quarum initium si insuper in puncto A ipso
collocatur , manifesto expressio coordinatarum z adipiscitur formam talem
ubi Q erit ordinis altioris quam secundi. Mutando dein situm axium ipsarum
X, y angulo M tali ut habeatur
tang2M=^^„
facile perspicitur, prodituram esse aequationem huius formae
z^\Txx-^\Vyy-\-^
quo pacto etiam tertiae conditioni satisfactum est. Quibus ita factis , patet
I. Si superficies curva secetur piano ipsi normali et per axem coordinata-
rum X t ranseunte , oriri curvam planam , cuius radius curvaturae in puncto A
fiat = ^ , signo positivo vel negativo indicante concavitatem vel convexitatem
versus plagam eam , versus quam coordinatae z sunt positivae.
II. Simili modo y erit in puncto A radius curvaturae curvae planae, quae
oritur per sectionem superficiei curvae cum piano per axes ipsarum y, z transeunte.

CIRCA SUPERFICIES CURVAS. 231
III. Statuendo ^ = rcoscp, j/ = rsincp, fit
z = 4-(rcoscp^H-Fsincp>r4-Q
linde colligitur, si Sectio fiat per planum superficiei in A normale et cum axe
ipsarum cc angulum cp efficiens, oriri curvam planam, cuius radius curvaturae in
puncto A sit
Tcoscp*+ Fsincp*
IV. Quoties itaque habetur T = V, radii curvaturae in cunctis planis
normalibus aequales erunt. Si vero T et F sunt inaequales , manifestum est,
quum Tcos^^-{- Fsincp^ pro quovis valore anguli cp cadat intra Tet F, radios
curvaturae in sectionibus principalibus , i n I et II consideratis , referri ad cur-
vaturas extremas , puta alterum ad curvaturam maximam , alterum ad minimam,
si Tet F eodem signo afFectae sint, contra alterum ad maximam convexitatem,
alterum ad maximam concavitatem, si T et F signis oppositis gaudeant. Hae
conclusiones omnia fere continent, quae ill. Euleb de curvatura superficierum cur-
varum primus docuit.
V. Mensura curvaturae superficiei curvae in puncto A autem nanciscitur
expressionem simplicissimam k = TV, unde habemus
Theorema. Mensura curvaturae in quovis superficiei puncto aequalis est
fractioni, cuius numerator unitas, denominator autem productum duorum radiorum
curvaturae extremorum in sectionibus per plana normalia.
Simul patet, mensuram curvaturae fieri positivam pro superficiebus concavo-
concavis vel convexo-convexis (quod discrimen non est essentiale), negativam vero
pro concavo-convexis. Si superficies constat e partibus utriusque generis , i n ea-
rum confiniis mensura curvaturae evanescens esse debebit. De indole superficie-
rum curvarum talium, in quibus mensura curvaturae ubique evanescit, infra plu-
ribus agetur.
9.
Formula generalis pro mensura curvaturae in fine art. 7 proposita, omnium
simplicissima est, quippe quae quinque tantum elementa implicat ; ad magis com-
plicatam, scilicet novem elementa involventem , d eferimur, si adhibere volumus

232 DISQÜISITIONES GENERALES
moduin primum indolem superficiei curvae exprimendi. Retinendo notatioiies
art. 4 insuper statuemus :
ddW _ p, ddPT __ ^, ddW _ „,
da:* — ^ ' dy* — ^ ' H^ — -"
d^^ _ p„ dd^T _ ^„ ^W^ __ „„
dy.dz -^ ' da;. dz ^' dx.dy —
ita ut tiat
dP = P'd^ + jR"dy + Q"d^
dQ = i2"djr+Q'dj/ + P"d;^
dE = Q"d^4-P"dy4-i2'd;^
p
lam quum habeatur ? = — ^ , i iivenimus per difFerentiationem
RRdt = — EdP+ PdJR = (P Q'—RP) d o; + (PP"— Pi2")dy + iPR'—RQ")dz
sive, eliminata d-s; adiumento aequationis Pd,2?+ Qd^^-j-Rds; = 0 ,
R'dt = {—RRr-]-2PRQ"—PPR)da^-\-{PRP"-J^QRQ"-PQR—IlRR')di/
Prorsus simili modo obtinemus
R'du = {PRP"-\- QR Q'-PQE- RRR")da^-\-{—RRQ'-{- '2QRP"-QQR')dy
Hinc itaque colligimus •
R^T= —RRP'-^2PRQ"—PPR'
R'U= PRP"^ QRQ"—PQR'—RRR"
R^V = —RRQ'-\-2QRP"—QQR'
Substituendo hos valores in formula art. 7 , obtinemus pro mensura curvaturae k
expressionem symmetricam sequentem :
{PP^QQ-{-RRfk
= PP[Q'R'—P'P")-^QQ{P'R'— Q"Q")^RR{P'Q'~R"R")
+ 2 QR{Q"R"--P'P")-Jr2PR{P"R"— QQ')-^2PQ{P"Q'—R'R")
10.
Formulam adhuc magis complicatam, puta e quindecim elementis con-
Üatam, obtinemus, si methodum generalem secundam, indolem superiicierum

CIRCA SUPERFICIES CURVAS. 233
curvarum exprimendi , sequimur. Magni tarnen momenti est , h anc quoque ela-
borare. Retinendo signa art. 4 , i nsuper statuemus
dda; ddx r ddx „
"d^ * ' d^d^ ' d^ ^
ddy __ ß ddy __ ^, ddy __ g"
dp^ ' d^.dy ' dj*
ddz ddz f ddz „ f
dP" ')' ' d^d^ T ' d^ T
Praeterea brevitatis caussa faciemus
bc — cb' = A
cd — ac = B
ah' — bd = C
Primo observamus, haberi Adx-\-Bdi/-\-Cdz =^ 0, sive d^; = — ~daj — ^dy;
quatenus itaque z spectatur tamquam functio ipsarum o?, ^, fit
dz A
dx— ^ — ~C
dz^ S
dy C
Porro deducimus , ex dx =^ adp-{-ddq, di/ = bdp-\-b'dq,
Cdp = b'dx — ddi/
Cdq = — bdaj-{-adi/
Hinc obtinemus differentialia completa ipsarum t, u
CM«=(B^^_C|^){R^-a'dy) + (C^-ß^^){6d..-«dy)
lam si in bis formulis siibstituimus
d^ ' , ' ' ,
— = ay4-ca— ay — ca
dB ,,.„„, ,
^=ay + ca— ay— ca
■^ z=ib'a-\-a'6' — ba! — a'ö
-— = b'a-\- d6 " — b a — d6 ' 38

atque perpendimus , valores difFerentialium dt, du sie prodeuntium, aequales
esse debere, independenter a differentialibus d<2?,dj/, quantitatibus Tda!-{-Udi/,
UdoJ-^ Vdi/ resp. inveniemus, post quasdam transformationes satis obvias:
C'T== aA b'b'-j- ^ B Vh'-\- 7 C h'h'
— 2a Abb'— 2n'Bbb'—2y'Cbb'
-\-aAbb-\--6"Bbb-\-i'Cbb
C^U^ —aAab'—-6Bab'—yCdb'
-^a'A{ab'-{-bd)-\-ß'B{ab'-j-ba)^iC{ab'-\-ba)
— a"Aab — ^"Bab — y'Cab
C^V= aA dd-\- € B äd-\- 7 C dd
— 2dAad — 2^'Bad — 2^Cad
-\-a'Aaa-{-W'Baa-\-fCaa
Si itaque brevitatis caussa statuimus
Aa-^Bß 4-C7 =D (1)
Aa-\-Bß'^Cj =D' . (2)
Aa"-{-B^"-]-CY = D" (3)
fit
C^T= Db'b'—2D'bb'-{-D"bh
C^ü=—Ddb'-Jr D\ab'-\-bd) — D"ab
C^V= Ddd—2D'ad-\-D"aa
Hinc invenimus , evolutione facta ,
C\TV- Uü) = (DD"—B'D'){ab'—bdf = [DD"—D'D')CC
et proin formulam pro mensura curvaturae
h —
DD"—D'D'
Formulae modo inventae iam aliam superstruemus , quae inter fertilissima
theoremata in doctrina de superficiebus curvis referenda est. Introducamus se-
quentes notationes :

CIRCA SUPERFICIES CUR VAS. 235
aa-\-bb -^cc = E
aa-\-bh' -{-cc = F
aa-\-b'b' -\-cc = G
aa -\-b^ +^7 = m (4)
aa' -\-bß' -\-cy' = m (5)
aa"-{-bß"-{-cy" = m" (6)
a'a -{-b'ß -\-cy = n (7)
aa'-\-b'ß' -\-c'Y = n (8)
da^b'ß"-{-cY ^n" (9)
ÄA-i-BB-^CC=EG — FF= A
Eliminemus ex aequationibus 1, 4, 7, quantitates ^, y, quod fit multipli-
cando illas per bc — cb\ b'C — c'B, cB — bC, etaddendo: ita oritur
(A{bc — cb')-{-a[b'C — cB)-\-a\cB — bC))a
= D[bc — cb')^m{b'C—cB)-\-n{cB — bC)
quam aequationem facile transformamus in hanc :
AD = a^-\-a{nF—m G)-{-a{mF~nE)
Simili modo eliminatio quantitatum a, y vel a, ö ex iisdem aequationibus sup-
peditat
BD = ^/l-]-b{nF—mG)-{-b'{mF—nE)
CD = y^-{-c[nF—mG)-\-c{mF—nE)
Multiplicando has tres aequationes per a!' , ö ", y" et addendo obtinemus
DD" := [a^"^m"-^yi')^^m(7iF—mG)-\-n{mF—nE) .... (10)
Si perinde tractamus aequationes 2, 5, 8 , prodit
AD' = a^-{-a{nF—mG)-\-a{mF—nE)
BD' = ^'ä-\-b[nF—rnG)^b'(mF—7iE)
CD' =^ y'^^c{nF — m'G)-\-c{m'F — n'E)
quibus aequationibus per a, ö', y' m ultiplicatis , additio suppeditat:
DD' =. {aa-i-^'ß'-\-Yy')^-\-m'(n'F~m'G)-{'n{m'F—nE)

Combinatio huius aequationis cum aequatione (10) producit
DD"— DD' = {aa-\-^ß"-{-yy"—aa'—■6'^'—iY)^
+ E {n'n— nn")-\-F{n m" — 2 mn'-\- mn")-\-G {mm — m m" )
lam patet esse ^ dq =z In
dF ^ dF ^ , dF , , dF „ , f dG ^ > d6
~=2m, -T~--27n, -T-z=m-\-n, j- =
d^ dq 'dp ' ' dy
sive
dJ=^ , d 6?
^ dF r , dF
dj ^ dp
dF . dF , 1 d G
Porro facile confirmatur, haberi
, d G
--i
dn' d rn" d m'
aa"+Ö6"+7Y"— aV— Ö'6'— yY = ^- ~ d p dp dq
dd^_, ddP , d dG
' d q^ ~^ dp.dq ^ • d p'
Quodsi iam has expressiones diversas in formula pro mensura curvaturae in fine
art. praec. eruta substituimus, pervenimus ad formulam sequentem, e solis quan-
titatibus F, F, G atque earum quotientibus differentialibus primi et secundi or-
dinis concinnatam:
A{EG-FFYk=:^ E(^-^-^^-^-^{^))
' ^dp ' d q dq 'dp dq ' d q •" dp ' d q dp dp '
, ^/d^ dG dF dFj,(^)\
~^^^dp ■dp^'^'dp ' d q'^^dq^ f
12.
Quum indefinite habeatur
da?2-|-d/ + d;^^ = £d/+2i^d;?. d^+ (7d^*
patet, \J{Ed27'-]-2Fdp.dq-\- Gdq^) esse expressionem generalem elementi li-
nearis in s uperficie curva. Docet itaque analysis in art. praec. explicata , ad in-
veniendam mensuram curvaturae haud opus esse formulis finitis , quae coordina-

CIRCA SUPERFICIES CURVAS. 237
tas X, y, z tamquam funetiones indeterminatarum jö, q exhibeaiit, sed sufficerc
expressionem generalem pro magnitudine cuiusvis elementi linearis. Progredia-
mur ad aliquot applicationes huius gravissimi theorematis.
Supponamiis , superficiem nostram curvam explicari posse in aliam supei*fi-
ciem , curvam seu planam , i ta ut cuivis puncto prioris superficiei per coordina-
tas X, y, z determinato respondeat punctum determinatum superficiei posterio-
ris , c uius coordinatae sint x, y\ z. Manifeste itaque x, y, z quoque conside-
rari possunt tamquam funetiones indeterminatarum p, q, unde pro elemento
\j{(}ix'' ~\-dy''^-\-^^'^) prodibit expressio talis
sl[E'dif^lF'dip.diq^G'dq')
denotantibus etiam E\ F', G' funetiones ipsarum p, q. At per ipsam notionem
expUcationis superficiei in superficiem patet , elementa in utraque superficie cor-
respondentia necessario aequalia esse, adeoque identice fieri
E = E', F=F\ G=^ G'
Formula itaque art. praec. sponte perducit ad egregium
Theorema. Si superficies curva in quamcunque aliam superficiem explicatur,
mensura curvaturae in singulis punctis invariata manet.
Manifeste quoque quaevis pars finita superficiei curvae post explicationem in
aliam superficiem eandem ciirvaturam integram retinebit.
Casum specialem, ad quem geometrae hactenus investigationes suas restrin-
xerunt, sistunt superficies in planum explicabiles. Theoria nostra sponte docet,
talium superficierum mensuram curvaturae in quo vis puncto fieri = 0, quocirca,
si earum indoles secundum modum tertium exprimitur, ubique erit
ddz ddz / ddz -^
quod criterium , dudum quidem notum , plerumque nostro saltem iudicio haud eo
rigore qui desiderari posset demonstratur.
13.
Quae in art. praec. exposuimus , cohaerent cum modo peculiari superficies
considerandi , summopere digno , qui a geometris diligenter excolatur. Scilicet
quatenus superficies consideratur nou tamquam limes solidi, sed tamquam soll-

dum, cuius dimensio una pro evanescente habetur, flexile quidem, sed non ex-
tensibile, qualitates superficiei partim a forma pendent, in quam illa reducta con-
cipitur, partim absolutae sunt, atque invariatae manent, in quamcunque formam
illa fiectatur. Ad has posteriores, quarum investigatio campum geometriae novum
fertilemque aperit, referendae sunt mensura curvaturae atque curvatura integra
eo sensu, quo liae expressiones a nobis accipiuntur; porro huc pertinet doctrina
de lineis brevissimis , pluraque alia , d e quibus in posterum agere nobis reserva-
mus. In hoc considerationis modo superficies plana atque superficies in planum
explicabilis, e. g. cylindrica, conica etc. tamquam essentialiter identicae spectan-
tur, modusque genuinus indolem superficiei ita consideratae generaliter expri-
mendi semper innititur formulae \Ji Edp^-\-2Fdp.dq-{- Gdq^) , quae nexum
elementi cum duabus indeterminatis p, q sistit. Sed antequam hoc argumentum
ulterius prosequamur, principia theoriae linearum brevissimarum in superficie
curva data praemittere oportet.
14.
Indoles lineae curvae in spatio generaliter ita datur , ut coordinatae sc, y, z
singulis illius punctis respondentes exhibeantur in forma functionum unius varia-
bilis , quam per w denotabimus. Longitudo talis lineae a puncto initiali arbi-
trario usque ad punctum, cuius coordinatae sunt x,y, z, exprimitur per integrale
^»•v((^-£)+(^:)+(t^f)
Si supponimus, situm lineae curvae variationem infinite parvam pati, ita ut coor-
dinatae singulorum punctorum accipiant variationes So?, hy, hz , variatio totius
longitudinis invenitur
~ J ^d a;. doa7;( d +^ ^ d + y. d ^ d Ä o + y d - ^ )- ) d z.dS^
quam expressionem in hanc formam transmutamus :
Ax .^x^ Ay .hy -\- A z.hz
v/(dx^ + dy^ + ds^r"
~ i ( ^' ^ •^ ^7^d^M:dy+Tz^) + ^^ •^ 7|d^^
In casu eo, ubi linea est brevissima inter puncta sua extrema, constat, ea, quae
hie sub signo integrali sunt, evanescere debere. Quatenus linea esse debet in su-

CIRCA SUPERFICIES CÜRVAS. 239
perficie data, cuius indoles exprimitur per aequationem Pda;-\- Qd^-^-Rdz = 0,
etiam variationes ^o), ^j/, ^ z satisfacere debent aequationi P5^-|- Q^i/-\- RBz = 0,
unde per principia nota facile colligitur , difFerentialia
A ^^ ^ dy ^ dz
s/idx' + dtf + dz'')' ^^{dx' + di/' + dz') ' \/{dx^ + di/^ + dz')
resp. quantitatibus P, Q, R proportionalia esse debere. lam sit dr elementum
lineae curvae , X punctum in superficie sphaerica repraesentans directionem huius
elementi, L punctum in superficie sphaerica repraesentans directionem normalis
in superficiem curvam ; denique sint S, t j, C coordinatae puncti X, atque X, Y, Z
coordinatae puncti L respectu centri sphaerae. Ita erit
da? = ^dr, dy = y\dr , dz = Cdr
unde colligimus, differentialia illa fieri d^, dr], dC. Et quum quantitates P. Q, R
proportionales sint ipsis X, Y, Z, character lineae brevissimae consistit in ae-
quationibus
dXj dFr) dYC
Ceterum facile perspicitur, \/( d ?~ -]- d tj^ + d C^) aequari arculo in superficie sphae-
rica, qui m ensurat angulum inter directiones tangentium in initio et fine elementi
dr, adeoque esse = -- , si p denotet radium curvaturae in hoc loco curvae bre-
vissimaei ; ta fiet
pd^ = Xdr, pdT] = Ydr, pdC = Zdr
15.
Supponamus , i n superficie curva a puncto dato Ä proficisci innumeras cur-
vas brevissimas , quas inter se distinguemus per angulum , quem constituit Sin-
gular um elementum primum cum elemento primo unius ex his lineis pro prima
assumtae: sit cp i lle angulus, vel generalius functio illius anguli, nee non r lon-
gitudo talis lineae brevissimae a puncto A usque ad punctum . cuius coordinatae
sunt 0?, 1/, z. Quum itaque valoribus determinatis variabilium r, z> respondeant
puncta determinata superficiei, coordinatae x,y, z considerari possunt tamquam
functiones ipsarum r, cp. Notationes X, L, $, t], C , X, F, Z in eadem significa-

tione retinebimus , i n qua in art. praec. acceptae fuenmt, modo indefinite ad
punctum indefinitum cuiuslibet linearum brevissimarum referantur.
Lineae brevissimae omnes , quae sunt aequalis longitudinis r, terminabun-
tur ad aliam lineam , cuius longitudinem ab initio arbitrario numeratam denota-
mus per v. Considerari poterit itaque v tamquam functio indeterminatarum r, cp,
et si per X' designamus punctum in superficie sphaerica respondens directioni ele-
menti dv, nee non per ^', t]', C coordinatas huius puncti respectu centri sphae-
rae , h abebimus :
da; __ 5./ d v d^ , du dz^ ^r dv
düä ''die' d^ ^•d9' dcp '"dcp
Hinc et
ex
dr " dr '»' dr ^
sequitur
dcp
da:
da; , d y dy d_z dz _ f-C I ^^'_i_r^") ^-1 = cos XX' -
dr ^dr dcp ' dr'dcp — U'n-TjT^i-T^^' i- dcp — ^^^'^'^-dcp
Membrum primum huius aequationis, quod etiam erit functio ipsarum r, cp, per
S denotamus; cuius difFerentiatio secundum r suppeditat:
dS _ ddx d^jddy ^|ddz_d£,^ ^Ud7^ "I~ld7^ ~rld^) )
dr dr* *dcp "" dr'' ' d cp " T" dr"* ' d cp "T"'2"- dcp
d r ' d cp " • d r d cp " • d r * d cp ' " ■^ ' d cp
Sed ^^-{-'Yi'Yi-\-C,C, = 1, adeoque ipsius differentiale =0; et per art. praec.
habemus, si etiam hie p denotat radium curvaturae in linea r,
dj X dY] r dC z
dr p ' dr p ' dr p
Ita obtinemus
^d r = ^p(^xe ' +F'T ] ''+'zc/')d.c^ p = -p.co sx x'.dc^p =:o "
quoniam manifesto X' iacet in circulo maximo, cuius polus L. Hinc itaque con-
cludimus, 8 independentem esse ab r et proin functionem solius cp . At pro r=0
manifesto fit v =z 0, et proin etiam ~ = 0, nee non 8=0, independenter a cp.
Necessario itaque generaliter esse debebit 8=0, adeoque cosXX' = 0, i.e.
XX'= 90*^. Hinc colligimus


\end{document}
